\documentclass[11pt]{article}
\usepackage{fontspec}
\setmainfont{Times New Roman}
\setsansfont{Arial}
\setmonofont{Consolas}
\usepackage{amsmath,amssymb,mathtools}
\usepackage{geometry}
\usepackage{microtype}
\geometry{a4paper,margin=2.28cm}
\setlength{\parindent}{1.15em}
\setlength{\parskip}{0pt}
\makeatletter
\renewcommand\footnotesize{\@setfontsize\footnotesize{7.5}{8.5}\selectfont}
\makeatother
\providecommand{\mo}{\mathfrak{o}}
\providecommand{\mM}{\mathfrak{M}}
\providecommand{\mU}{\mathfrak{U}}
\providecommand{\mA}{\mathfrak{A}}
\providecommand{\mB}{\mathfrak{B}}
\emergencystretch=120em
\allowdisplaybreaks
\sloppy
\hbadness=10000
\vbadness=10000
\hfuzz=2pt
\vfuzz=10pt
\raggedbottom
\begin{document}
% Унит: Папер34 сецтион16 в001. Дирецт Интерславиц Латин ауториты транслатион фром аудитед Герман соурце линес 1247--1310, сегментс P34-С0212--P34-С0219; P34-С0220/Сецтион17 ресервед фор некст транцхе. Зенодо рецорд 20822156 цхецкед фор тис унит ат 2026-06-25Т01:59З УТЦ вит но филе кеы/сизе/цхецксум цханге агаинст the Сецтион15 префлигхт.
\section*{34. Хиперкомплексне величины и теорија представјень}
\emph{Продолженье и завршенье: глава III, §§15--19, и глава IV, §§20--26.}

\section*{III. поглавје. Теорија модулов и представјень}

\subsection*{§ 16. Редуцибилне представјеньа}

Ако $\mU$ јест подмодул модула представјеньа $\mM$, и ако можно јест выбрати базу за $\mM$, составльену из базы $z_1,\ldots,z_t$ од $\mU$, дополньену с $y_1,\ldots,y_r$, то јест
\[
        \mM=y_1K+\cdots+y_rK+z_1K+\cdots+z_tK,\textsuperscript{16)}
\]
тогда представјеньа имајут форму
\[
        C=\begin{pmatrix}R&0\\ S&T\end{pmatrix},
\]
где матрице $T$ саме по себе творјат представјенье $t$-го степена, породжено модулом $\mU$, а матрице $R$ творјат представјенье $r$-го степена, породжено модулом $\mM/\mU$.

Доказ. Јест
\[
        (cy_1,\ldots,cy_r,cz_1,\ldots,cz_t)
        =(y_1,\ldots,y_r,z_1,\ldots,z_t)C.
\]
Понеже $cz_i$, како елементы од $\mU$, изразјајут се само чрез $z_i$, в матрици $C$ вправо горе стоји нула. Ако остале делы од $C$ в показаној порјадности назовемо $R,S,T$, тогда особливо
\[
        (cz_1,\ldots,cz_t)=(z_1,\ldots,z_t)T;
\]
зато $T$ творјат представјенье, посредовано чрез $\mU$. Далье
\[
        (cy_1,\ldots,cy_r)\equiv (y_1,\ldots,y_r)R\pmod{\mU},
\]
меджутым $y_i$ творјат базу, линеарно независну модуло $\mU$. Зато $R$ творјат представјенье, породжено модулом $\mM/\mU$.

Обратно, ако дано јест ``редуцибилно'' представјенье
\[
        C=\begin{pmatrix}R&0\\ S&T\end{pmatrix},
\]
где $R$ и $T$ сут квадратне матрице, тогда в приналежном модулу представјеньа продукты всакого $c$ с последними $t$ базовыми елементами $z_1,\ldots,z_t$ изразјајут се само чрез ньих, т. ј. $\mU=(z_1,\ldots,z_t)$ јест подмодул.

{\footnotesize\noindent 16) Иными словами: когда $\mU$ како $K$-модул јест директны сумманд. Ако $K$ јест тело, то предпоставјенье всегда јест изполньено.\par}

Последок. Нехај
\[
        \mM\supset \mU_1\supset \mU_2\supset\cdots\supset \mU_e=0
\]
буде композицијны ред за $\mM$, и нехај всакы $\mU_i$ в предходном буде директны сумманд како $K$-модул, тако же можно јест выбрати базу
\[
        z_{11},\ldots,z_{1r_1},\ z_{21},\ldots,z_{2r_2},\ldots,
        z_{e1},\ldots,z_{er_e}
\]
тако же $z_{ik}$ с $i>\nu$ творјат базу за $\mU_\nu$. Тогда представјеньа, посредована чрез $\mM$, имајут форму
\[
 C=\begin{pmatrix}
 R_{11}&0&\cdots&0\\
 R_{12}&R_{22}&\cdots&0\\
 \vdots&\vdots&\ddots&\vdots\\
 R_{1e}&R_{2e}&\cdots&R_{ee}
 \end{pmatrix},                                      \tag{2}
\]
и $R_{\nu\nu}$ творјат представјеньа, посредована чрез композицијне факторы $\mU_{\nu-1}/\mU_\nu$.

Јединачне представјеньа $R_{\nu\nu}$ сут иредуцибилне, понеже композицијне факторы $\mU_{\nu-1}/\mU_\nu$ сут просте. Ако предпоставити, же $K$ јест тело, тогда и обратно всако представјенье, најдалье редуцирано до формы (2), води к композицијному реду. По теорему Јордана--Холдера композицијне факторы сут једнозначне до операторного изоморфизма; зато $R_{\nu\nu}$ сут једнозначно определьене до еквивалентных представјень и до порјадности.

Ако $\mM=\mA+\mB$ јест директна сума двух модулов представјеньа $\mA=(y_1,\ldots,y_r)$, $\mB=(z_1,\ldots,z_s)$, тогда представјенье, посредовано чрез базу $(y_1,\ldots,y_r,z_1,\ldots,z_s)$, јавно имаје форму
\[
        C=\begin{pmatrix}R&0\\0&S\end{pmatrix},
\]
где $R$ означује представјенье елемента $c$, посредовано чрез $\mA$, а $S$ оно, посредовано чрез $\mB$, и обратно. Ако модул $\mM$ најпрво представити како директну суму директно неразложимых составных честиј и за ньих изискати композицијне реды како выше, тогда при погодном выбору базы представјенье имаје одповедајучу блокову форму. Ако знову $K$ јест тело, тогда и обратно всако најдалье редуцирано представјенье води к представјеньу модула чрез директно неразложиме сумманды и композицијне факторы в ньих. По теорему Крулла и Отта Сцхмидта класы директно неразложимых составных честиј представјеньа сут једнозначно определьене до порјадности.

Ако модул $\mM$ јест пополно редуцибилны, тогда вси блокы состојат из јединого иредуцибилного представјеньа, и представјенье называ се пополно редуцибилным.

\end{document}

